\documentclass[11pt]{article}
\usepackage[a4paper,margin=1in]{geometry}
\usepackage{amsmath,amssymb}
\usepackage{authblk}
\title{Bell-Certified Randomness and Asymptotic Pattern Completeness: The AST--Q Framework}
\author{Nilansh Anand\\(conceptual origin)\\Collaborating AI Systems\\(Perplexity, ChatGPT, Gemini)}
\date{\today}
\begin{document}
\maketitle

\begin{abstract}
Bell inequality violations show that certain quantum processes produce outcomes incompatible with any local deterministic hidden-variable model, enabling device-independent certification of intrinsic randomness via Bell-type experiments. Classical probability theory implies that sufficiently long sequences of independent random outcomes almost surely realize every finite pattern as a subsequence. This note combines these results into the AST--Q (Anand--Stochastic Totality--Quantum) framework: if a physical system generates Bell-certified random bits over unbounded time, then its output asymptotically exhibits pattern completeness, while any finite observer necessarily experiences that totality only through stochastic samples.
\end{abstract}

\section{CHSH inequality in local hidden-variable theories}

Consider a standard CHSH scenario with two distant observers, Alice and Bob. Alice chooses a setting $a$ or $a'$, Bob chooses a setting $b$ or $b'$, and outcomes are $A(a,\lambda), A(a',\lambda), B(b,\lambda), B(b',\lambda) \in \{-1,+1\}$, where $\lambda$ denotes hidden variables.

A \emph{local hidden-variable} (LHV) model assumes:
\begin{itemize}
  \item \textbf{Realism:} For each $\lambda$, the four outcomes $A(a,\lambda), A(a',\lambda), B(b,\lambda), B(b',\lambda)$ are well-defined, independent of which settings are actually chosen.
  \item \textbf{Locality:} Alice's outcome depends only on $(a,\lambda)$, Bob's only on $(b,\lambda)$; there is no dependence on the distant setting.
  \item \textbf{Statistical distribution:} There is a probability density $\rho(\lambda) \ge 0$ with $\int \rho(\lambda)\,d\lambda = 1$.
\end{itemize}

Define the correlation for settings $x \in \{a,a'\}$, $y \in \{b,b'\}$ as
\begin{equation}
E(x,y) = \int d\lambda\, \rho(\lambda)\, A(x,\lambda) B(y,\lambda).
\end{equation}

The CHSH combination is
\begin{equation}
S = E(a,b) + E(a,b') + E(a',b) - E(a',b').
\end{equation}

For each hidden variable value $\lambda$, define the ``single-run'' quantity
\begin{equation}
C(\lambda) = A(a,\lambda)B(b,\lambda)
+ A(a,\lambda)B(b',\lambda)
+ A(a',\lambda)B(b,\lambda)
- A(a',\lambda)B(b',\lambda).
\end{equation}
Then $S = \int d\lambda\,\rho(\lambda)\,C(\lambda)$.

Factor $C(\lambda)$ as
\begin{equation}
C(\lambda) = A(a,\lambda)\big[B(b,\lambda)+B(b',\lambda)\big]
+ A(a',\lambda)\big[B(b,\lambda)-B(b',\lambda)\big].
\end{equation}

Since each $B(\cdot,\lambda) \in \{-1,+1\}$, the sums $B(b,\lambda)+B(b',\lambda)$ and $B(b,\lambda)-B(b',\lambda)$ each lie in $\{-2,0,2\}$, and for any fixed $\lambda$ exactly one of these two brackets is nonzero and equal to $\pm 2$. The nonzero term is multiplied by some $A(\cdot,\lambda)\in\{-1,+1\}$, so
\begin{equation}
C(\lambda) \in \{-2,+2\} \quad \Rightarrow \quad |C(\lambda)| = 2.
\end{equation}

Using $S = \int d\lambda\,\rho(\lambda)C(\lambda)$ and the triangle inequality,
\begin{equation}
|S| = \left|\int d\lambda\,\rho(\lambda)C(\lambda)\right|
\le \int d\lambda\,\rho(\lambda)|C(\lambda)|
\le \int d\lambda\,\rho(\lambda)\,2 = 2.
\end{equation}

Thus any LHV model obeys the CHSH inequality
\begin{equation}
|S| \le 2.
\end{equation}
This constitutes the \emph{classical floor} for correlations under locality and realism.

\section{Quantum violation and Tsirelson's bound}

Now consider two spin-$1/2$ particles in the singlet state
\begin{equation}
|\psi\rangle = \frac{1}{\sqrt{2}}\Big(|\uparrow\rangle_A|\downarrow\rangle_B - |\downarrow\rangle_A|\uparrow\rangle_B\Big).
\end{equation}

Let Alice and Bob measure spin along unit vectors $\mathbf{a}, \mathbf{a}', \mathbf{b}, \mathbf{b}'$ in a common plane. The observables are $\hat{A} = \boldsymbol{\sigma}_A\cdot\mathbf{a}$, $\hat{A}' = \boldsymbol{\sigma}_A\cdot\mathbf{a}'$, $\hat{B} = \boldsymbol{\sigma}_B\cdot\mathbf{b}$, $\hat{B}' = \boldsymbol{\sigma}_B\cdot\mathbf{b}'$, with eigenvalues $\pm 1$. For the singlet state, the quantum correlation is
\begin{equation}
E(\mathbf{a},\mathbf{b}) = \langle \psi | \hat{A}\otimes \hat{B} | \psi \rangle = -\,\mathbf{a}\cdot\mathbf{b} = -\cos\theta,
\end{equation}
where $\theta$ is the angle between $\mathbf{a}$ and $\mathbf{b}$.

Choose measurement directions that maximize CHSH violation, for example
\begin{equation}
\mathbf{a} = (1,0), \quad 
\mathbf{a}' = (0,1), \quad
\mathbf{b} = \tfrac{1}{\sqrt{2}}(1,1), \quad
\mathbf{b}' = \tfrac{1}{\sqrt{2}}(1,-1).
\end{equation}
Then
\begin{align}
E(a,b)   &= -\mathbf{a}\cdot\mathbf{b}   = -\tfrac{1}{\sqrt{2}},\\
E(a,b')  &= -\mathbf{a}\cdot\mathbf{b}'  = -\tfrac{1}{\sqrt{2}},\\
E(a',b)  &= -\mathbf{a}'\cdot\mathbf{b}  = -\tfrac{1}{\sqrt{2}},\\
E(a',b') &= -\mathbf{a}'\cdot\mathbf{b}' = +\tfrac{1}{\sqrt{2}}.
\end{align}

Using the standard CHSH combination
\begin{equation}
S_{\text{QM}} = E(a,b) + E(a,b') + E(a',b) - E(a',b'),
\end{equation}
we obtain
\begin{equation}
S_{\text{QM}} = -\tfrac{1}{\sqrt{2}} - \tfrac{1}{\sqrt{2}} - \tfrac{1}{\sqrt{2}} - \tfrac{1}{\sqrt{2}} = -2\sqrt{2},
\end{equation}
so
\begin{equation}
|S_{\text{QM}}| = 2\sqrt{2} > 2.
\end{equation}

This is Tsirelson's bound: for quantum systems, $|S| \le 2\sqrt{2}$, and this upper limit is attainable but cannot be exceeded. Experiments implementing CHSH-type tests with entangled photons, ions, and superconducting qubits consistently observe values of $|S|$ approaching this quantum limit and significantly exceeding 2, thereby ruling out all local hidden-variable models.

In the AST--Q language, this establishes a \emph{certified non-determinism engine}: there exist physical processes whose outcomes cannot be modeled as pre-assigned, local deterministic values and whose randomness can be certified in a device-independent manner.

\section{Infinite randomness and pattern completeness}

On the mathematical side, consider an idealized infinite sequence of independent, identically distributed random variables taking values in a finite alphabet with full support (for example, fair coin tosses). The \emph{infinite monkey theorem} and related Borel--Cantelli arguments imply that, with probability 1, such a sequence contains every finite word from the alphabet infinitely many times as a contiguous subsequence.

More precisely, let $\{X_n\}_{n\ge 1}$ be i.i.d. with $\mathbb{P}(X_n = x) > 0$ for all letters $x$. For any fixed finite word $w = (w_1,\dots,w_k)$, define the events ``$w$ starts at position $n$'' by
\begin{equation}
A_n = \{ X_n = w_1, X_{n+1} = w_2, \dots, X_{n+k-1} = w_k \}.
\end{equation}
These events are independent and have constant probability $p = \mathbb{P}(A_n) > 0$. Since $\sum_n \mathbb{P}(A_n) = \infty$, the second Borel--Cantelli lemma implies that almost surely infinitely many $A_n$ occur. Thus, for almost every realization of the process, every finite pattern is realized infinitely often as a subsequence.

This result holds \emph{almost surely} (with probability 1), though any finite prefix of the sequence need not exhibit full pattern coverage. In practice, ``totality'' is only visible in the infinite limit.

Algorithmic randomness refines this picture: an infinite binary sequence that is Martin--L\\"of random (incompressible in the sense of Kolmogorov complexity) also contains every finite binary word as a contiguous block. Such sequences serve as canonical examples of maximally random infinite objects.

\section{The AST--Q framework}

The AST--Q (Anand--Stochastic Totality--Quantum) framework is an interpretive layer connecting these two pillars:
\begin{enumerate}
  \item \textbf{Certified non-determinism (Engine):} Bell inequality violations and Tsirelson's bound show that certain quantum processes generate outputs that cannot be described by local deterministic hidden variables, and that these outputs can be used in device-independent randomness-generation protocols that quantify the min-entropy of the produced bits.

  \item \textbf{Asymptotic pattern completeness (Fuel):} Classical probability and algorithmic randomness imply that an ideal sequence of independent random bits, extended without bound, almost surely realizes every finite binary pattern as a subsequence, and that maximally random sequences encode irreducible information.

  \item \textbf{Finite observers (Vehicle):} Any physical observer (human, computer, detector) can access only a finite prefix of the total output of such a process. From this limited vantage point, the underlying pattern-complete structure is visible only as stochastic behavior: randomness is the low-resolution interface of an effectively infinite configuration space to a finite observer.
\end{enumerate}

A concise AST--Q statement is:

\medskip
\noindent
\emph{Bell-certified randomness provides a physically grounded stochastic source whose unbounded extension asymptotically exhibits pattern completeness, while any finite observer accesses only finite stochastic projections of this totality.}

\medskip
This statement does not assert that the universe is literally a ``Library of Babel,'' nor that Bell tests prove cosmological infinity or any specific interpretation of quantum mechanics. It instead identifies a controlled bridge: Bell-certified randomness furnishes the physical engine, while classical and algorithmic randomness theory describe the asymptotic informational behavior.

\end{document}
